%========================================
% Section：[8] Saṃvegaparikittanapāṭha
%========================================
\section*{m8 | Saṃvegaparikittanapāṭha}

% ---------- 第一段 ----------
\begin{chantsubsection}
  \chantlineinline
    {Idha tathāgato loke uppanno arahaṃ sammāsambuddho,}
    {Here in this world the Tathāgata, an arahant, perfectly awakened, has arisen.}
  \chantlineinline
    {Dhammo ca desito niyyāniko upasamiko parinibbāniko
     sambodhagāmī sugatappavedito.}
    {And the Dhamma has been taught that leads out, brings peace, leads to final Nibbāna,
     leads to awakening, proclaimed by the Well-gone One.}
  \chantlineinline
    {Mayantaṃ dhammaṃ sutvā evaṃ jānāma.}
    {Having heard that Dhamma, we understand thus:}
  \chantlineinline
    {Jātipi dukkhā, jarāpi dukkhā, maraṇampi dukkhaṃ,}
    {birth is suffering, aging is suffering, death is suffering;}
  \chantlineinline
    {sokaparidevadukkhadomanassupāyāsāpi dukkhā,}
    {sorrow, lamentation, pain, grief and despair are suffering;}
  \chantlineinline
    {Appiyehi sampayogo dukkho, piyehi vippayogo dukkho,}
    {association with the disliked is suffering, separation from the loved is suffering;}
  \chantlineinline
    {Yampicchaṃ na labhati, tampi dukkhaṃ.}
    {not getting what one wants is also suffering.}
\end{chantsubsection}
\chantgap
% ---------- 第二段 ----------
\begin{chantsubsection}
  \chantlineinline
    {Saṅkhittena pañcupādānakkhandhā dukkhā.}
    {In brief, the five aggregates subject to clinging are suffering.}
  \chantlineinline
    {Seyyathīdaṃ, rūpūpādānakkhandho vedanūpādānakkhandho
     saññūpādānakkhandho saṅkhārūpādānakkhandho
     viññāṇūpādānakkhandho,}
    {That is to say: the aggregate of clinging to form, to feeling, to perception,
     to formations, and to consciousness.}
  \chantlineinline
    {yesaṃ pariññāya dharamāno so bhagavā evaṃ bahulaṃ sāvake vineti.}
    {For the full understanding of these, while still alive, the Blessed One often
     trains his disciples in this way.}
  \chantlineinline
    {Evaṃbhāgā ca panassa bhagavato sāvakesu anusāsanī bahulā pavattati.}
    {And in this way his instruction among the disciples frequently proceeds.}
  \chantlineinline
    {Rūpaṃ aniccaṃ, vedanā aniccā, saññā aniccā, saṅkhārā aniccā,
     viññāṇaṃ aniccaṃ,}
    {Form is impermanent, feeling is impermanent, perception is impermanent,
     formations are impermanent, consciousness is impermanent;}
  \chantlineinline
    {rūpaṃ anattā, vedanā anattā, saññā anattā, saṅkhārā anattā,
     viññāṇaṃ anattā,}
    {form is not-self, feeling is not-self, perception is not-self,
     formations are not-self, consciousness is not-self;}
  \chantlineinline
    {sabbe saṅkhārā aniccā, sabbe dhammā anattāti.}
    {all conditioned things are impermanent, all dhammas are not-self.}
\end{chantsubsection}
\chantgap
% ---------- 第三段 ----------
\begin{chantsubsection}
  \chantlineinline
    {Te\footnote{For ladies, change \emph{Te} to \emph{Tā}.} mayaṃ, otiṇṇāmaha
     jātiyā jarāmaraṇena sokehi paridevehi dukkhehi domanassehi upāyāsehi.}
    {We, who are sunk in birth, aging and death, in sorrow, lamentation, pain,
     grief and despair,}
  \chantlineinline
    {Dukkhotiṇṇā dukkhaparetā, appevanāmimassa kevalassa
     dukkhakkhandhassa antakiriyā paññāyethāti.}
    {overwhelmed by suffering and oppressed by suffering, may the ending of this
     entire mass of suffering become apparent to us.}
  \chantlineinline
    {Ciraparinibbutampi taṃ bhagavantaṃ saraṇaṃ gatā,
     dhammañca saṅghañca,}
    {Though that Blessed One has long attained final Nibbāna, we have gone for
     refuge to that Blessed One, to the Dhamma and to the Saṅgha,}
  \chantlineinline
    {tassa bhagavato sāsanaṃ yathāsati yathābalaṃ manasikaroma
     anupaṭipajjāma.}
    {and we attend to that Blessed One’s teaching and practise it according to
     our mindfulness and strength.}
  \chantlineinline
    {Sā sā no paṭipatti imassa kevalassa dukkhakkhandassa
     antakiriyāya saṃvattatu.}
    {May this practice of ours lead to the ending of this entire mass of suffering.}
\end{chantsubsection}
